\section{Формализация критерия обнаружения регрессий и модель процесса профилирования в CI/CD}
\label{sec:razdel2}

\subsection{Модель процесса профилирования в CI/CD-пайплайне}
\label{sec:2.1}

Процесс обнаружения регрессии производительности рассматривается как
конечный автомат, привязанный к событию открытия или обновления pull
request (рис.~\ref{fig:2.1}). Ключевое проектное решение, вытекающее из
вывода раздела~\ref{sec:1.2}, — после запуска бенчмарков процесс
\emph{разветвляется} на два параллельных, статистически независимых шага:
сравнение агрегатов через \texttt{benchstat} (решающий шаг) и снятие
pprof-профилей с их сравнением через \texttt{diff\_base} (диагностический
шаг, не влияющий на исход сборки, но прикладываемый к отчёту).

\begin{figure}[htbp]
\centering
\begin{tikzpicture}[
  node distance=8mm and 12mm,
  box/.style={draw, rounded corners, minimum height=8mm, minimum width=28mm, align=center, font=\small},
  decision/.style={draw, diamond, aspect=2.2, minimum width=30mm, align=center, font=\small},
  arr/.style={-{Latex[length=2mm]}}
]
\node[box] (trigger) {PR открыт /\\обновлён};
\node[box, below=of trigger] (checkout) {checkout + setup Go\\+ install benchstat};
\node[box, below=of checkout] (bench) {go test -bench\\-count=N (эталон уже есть)};
\node[decision, below=10mm of bench] (split) {два\\параллельных\\шага};
\node[box, below left=10mm and -6mm of split] (compare) {benchstat vs.\\baseline.txt};
\node[box, below right=10mm and -6mm of split] (profile) {pprof-профиль vs.\\baseline-*.pprof};
\node[decision, below=16mm of compare] (dec) {$\Delta_m > T$ \&\&\\$p_m < \alpha$?};
\node[box, below=8mm of dec, fill=red!12] (fail) {FAIL: отчёт в PR,\\merge заблокирован};
\node[box, right=14mm of fail, fill=green!12] (ok) {OK: merge разрешён};
\node[box, below=8mm of profile] (artifact) {top-N функций в\\step summary + артефакт};

\draw[arr] (trigger) -- (checkout);
\draw[arr] (checkout) -- (bench);
\draw[arr] (bench) -- (split);
\draw[arr] (split) -- (compare);
\draw[arr] (split) -- (profile);
\draw[arr] (compare) -- (dec);
\draw[arr] (dec) -- node[left,font=\scriptsize]{да} (fail);
\draw[arr] (dec) -| node[above,pos=0.25,font=\scriptsize]{нет} (ok);
\draw[arr] (profile) -- (artifact);
\end{tikzpicture}
\caption{Модель процесса профилирования и обнаружения регрессии в CI/CD-пайплайне}
\label{fig:2.1}
\end{figure}

Существенное свойство модели: диагностический шаг (правая ветвь,
pprof-профиль) выполняется \texttt{always()} в терминах GitHub Actions —
то есть независимо от исхода решающего шага, включая случай его падения. Это
позволяет получить профиль-объяснение причины регрессии в том же прогоне
CI, в котором она была обнаружена, не дожидаясь отдельного ручного
профилирования разработчиком постфактум (что было отмечено как ограничение
подхода "ручное профилирование" в таблице~\ref{tab:1.1}).

\subsection{Формализация статистического критерия обнаружения регрессии}
\label{sec:2.2}

Пусть на эталонном коммите (main) и на кандидатном коммите (PR) для каждого
бенчмарка $b \in B$ и каждой метрики $m \in \{\text{ns/op}, \text{B/op},
\text{allocs/op}\}$ независимым запуском \texttt{go test -bench -count=$n$}
получены выборки
$$
X_{b,m} = \{x_1, \dots, x_n\}, \qquad Y_{b,m} = \{y_1, \dots, y_k\},
$$
где $X_{b,m}$ — измерения на эталоне, $Y_{b,m}$ — на кандидате
($n$ и $k$ обычно совпадают и равны параметру \texttt{-count}, но формально
не обязаны). Все три метрики в данной работе относятся к классу
"меньше — лучше" (lower-is-better), поэтому регрессией считается
\emph{рост} значения.

\texttt{benchstat} оценивает центральную тенденцию каждой выборки медианой
$\hat\mu(X_{b,m})$, $\hat\mu(X_{b,m})$ и вычисляет относительное изменение

\begin{equation}
\Delta_{b,m} = \frac{\hat\mu(Y_{b,m}) - \hat\mu(X_{b,m})}{\hat\mu(X_{b,m})} \times 100\%.
\label{eq:delta}
\end{equation}

Статистическая значимость отличия оценивается $p$-значением
$p_{b,m}$ критерия Манна --- Уитни (используется \texttt{benchstat} по
умолчанию как непараметрический критерий, устойчивый к отклонению
распределения задержек от нормального — задержки почти всегда
правосторонне скошены из-за редких выбросов планирования ОС) либо критерия
Уэлча~\cite{welch1947} для приблизительно нормальных выборок. Регрессия по
конкретной метрике конкретного бенчмарка объявляется тогда и только тогда,
когда одновременно выполнены статистическая и практическая значимость
отклонения:

\begin{equation}
\mathrm{regression}(b, m) \;\Longleftrightarrow\; \bigl(\Delta_{b,m} > T\bigr) \;\wedge\; \bigl(p_{b,m} < \alpha\bigr),
\label{eq:criterion}
\end{equation}

где $T$ — порог относительного отклонения в процентах (параметр
\texttt{THRESHOLD} в реализации, раздел~\ref{sec:razdel4}), $\alpha$ —
уровень значимости (параметр \texttt{ALPHA}). Условие $\Delta_{b,m} > T$
без учёта $p_{b,m}$ ловило бы шумовые колебания как ложные регрессии при
малом $T$; условие $p_{b,m} < \alpha$ без учёта $\Delta_{b,m}$ ловило бы
статистически значимые, но практически пренебрежимо малые отклонения
(например, 0,3\% при большой выборке) — сочетание обоих условий необходимо
именно для того, чтобы отделять шум от практически значимых регрессий; эта
гипотеза проверяется экспериментально в разделе~\ref{sec:5.2} перебором
сетки значений $T$ и $\alpha$.

Итоговое решение по всему прогону — дизъюнкция по всем бенчмаркам и
метрикам:

\begin{equation}
\mathrm{FAIL} \;\Longleftrightarrow\; \exists\, b \in B,\, m \in M : \mathrm{regression}(b, m).
\label{eq:fail}
\end{equation}

\textbf{Диагностический критерий на уровне профиля.} В отличие от
критерия~\eqref{eq:criterion}, сравнение pprof-профилей не участвует в
принятии решения FAIL/OK, а используется описательно: для CPU-профиля
каждая функция $f$ характеризуется числом сэмплов \texttt{flat} (время,
проведённое непосредственно в теле функции) и \texttt{cum} (суммарно с
учётом вызываемых функций); \texttt{go tool pprof -diff\_base} вычисляет

\begin{equation}
\delta_f = s_f^{\text{cur}} - s_f^{\text{base}}
\label{eq:profdiff}
\end{equation}

раздельно для \texttt{flat} и \texttt{cum}, где $s_f$ — число сэмплов,
приходящихся на функцию $f$. Функции ранжируются по $|\delta_f|$, и top-$N$
(в реализации $N=10$) выводятся в отчёт. Аналогичное выражение применяется
к heap-профилю с типом выборки \texttt{alloc\_space} (кумулятивный объём
выделенной памяти на функцию за время прогона).

\textbf{Связь с алгоритмической сложностью.} Отдельный частный случай
регрессии, воспроизведённый в разделе~\ref{sec:razdel5} (сценарий
\texttt{regression/\allowbreak quadratic-dedup}), — изменение класса временной
сложности алгоритма. Если исходная реализация имеет сложность $T(n) = O(n)$
(проход по хеш-таблице), а изменённая — $T(n) = O(n^2)$ (вложенный линейный
поиск, определение по~\cite{clrs}), то при фиксированном размере входа
$n=\text{const}$, используемом в бенчмарке, оба варианта дают лишь
\emph{константный} множитель времени выполнения — сам факт квадратичной
сложности не проявляется в единственном измерении при фиксированном $n$.
Изменение класса сложности обнаруживается критерием~\eqref{eq:criterion}
косвенно, через рост $\Delta_{b,m}$ при неизменном $n$, но не отличается
формально от "обычного" замедления в $k$ раз без изменения сложности — для
строгого различения потребовался бы бенчмарк с варьируемым $n$ и оценка
показателя роста, что выходит за рамки данной работы и отмечено как
направление дальнейшего исследования в заключении.

\subsection{Концептуальная схема интеграции двухуровневого профилирования в пайплайн}
\label{sec:2.3}

Формализованный процесс (рис.~\ref{fig:2.1}) и критерий~\eqref{eq:criterion}
определяют требования к артефактам, которые должна хранить и обновлять
система (рис.~\ref{fig:2.2}): эталонный агрегированный отчёт
(\texttt{baseline.txt}), эталонные pprof-профили
(\texttt{baseline-cpu.pprof}, \texttt{baseline-mem.pprof}) и параметры
критерия ($T$, $\alpha$, \texttt{-count}), являющиеся внешними по отношению
к коду параметрами запуска, а не константами внутри приложения — это
позволяет проводить эксперименты раздела~\ref{sec:razdel5} (перебор сетки
$T \times \alpha$) без изменения кода прототипа.

\begin{figure}[htbp]
\centering
\begin{tikzpicture}[
  node distance=10mm and 10mm,
  store/.style={draw, cylinder, shape border rotate=90, aspect=0.25, minimum height=14mm, minimum width=24mm, align=center, font=\small},
  proc/.style={draw, rounded corners, minimum height=9mm, minimum width=32mm, align=center, font=\small},
  arr/.style={-{Latex[length=2mm]}}
]
\node[store] (base) {baseline.txt};
\node[store, below=6mm of base] (basecpu) {baseline-*.pprof};
\node[proc, right=22mm of base] (app) {app/ (Go-приложение\\+ бенчмарки)};
\node[proc, below=of app] (compare) {compare.sh\\(benchstat + критерий)};
\node[proc, below=of compare] (diff) {profile-diff.sh\\(pprof diff\_base)};
\node[proc, right=22mm of app] (ci) {CI (GitHub Actions):\\perf-check.yml};
\node[proc, below=of ci] (summary) {\$GITHUB\_STEP\_SUMMARY\\+ pprof-artifact};

\draw[arr] (base) -- (compare);
\draw[arr] (basecpu) -- (diff);
\draw[arr] (app) -- (compare);
\draw[arr] (app) -- (diff);
\draw[arr] (ci) -- (compare);
\draw[arr] (ci) -- (diff);
\draw[arr] (compare) -- (summary);
\draw[arr] (diff) -- (summary);
\draw[arr] (ci) -- (summary);
\end{tikzpicture}
\caption{Концептуальная схема артефактов и компонентов интеграции профилирования в CI/CD}
\label{fig:2.2}
\end{figure}

\subsection{Выводы}

\begin{enumerate}
\item Разработана модель процесса профилирования в CI/CD-пайплайне
(рис.~\ref{fig:2.1}) в виде схемы с ветвлением на решающий шаг
(\texttt{benchstat}-сравнение) и диагностический шаг (pprof-\texttt{diff\_base}),
выполняемые параллельно и независимо в рамках одного прогона CI.
\item Формализован критерий обнаружения регрессии~\eqref{eq:criterion},
сочетающий практическую значимость отклонения $\Delta_{b,m}$
(выражение~\eqref{eq:delta}) и статистическую значимость $p_{b,m}$; итоговое
решение по сборке определено как дизъюнкция по всем бенчмаркам и метрикам
(выражение~\eqref{eq:fail}).
\item Формализован диагностический критерий на уровне отдельной функции
профиля~\eqref{eq:profdiff}, не влияющий на решение FAIL/OK, но
предоставляющий объяснение причины регрессии без отдельного ручного
профилирования.
\item Показано (со ссылкой на определение алгоритмической сложности
по~\cite{clrs}), что изменение класса временной сложности алгоритма при
фиксированном размере входа $n$ бенчмарка формально неотличимо от обычного
замедления в константное число раз — это ограничение модели явно
зафиксировано и учтено при интерпретации сценария
\texttt{regression/\allowbreak quadratic-dedup} в разделе~\ref{sec:razdel5}.
\item Спроектирована концептуальная схема артефактов (рис.~\ref{fig:2.2}):
параметры критерия ($T$, $\alpha$, \texttt{-count}) и эталонные данные
(текстовый отчёт и бинарные pprof-профили) вынесены за пределы кода
приложения, что позволяет проводить экспериментальную апробацию
(раздел~\ref{sec:razdel5}) без изменения исходного кода прототипа.
\end{enumerate}
